% Developing a Feasible Neuroplastic Mechanical Neuron in the Domain of Soft Robotics
% LaTeX version of WernerM-Article2.docx (converted 2026-10-01).
% Build: pdflatex main && bibtex main && pdflatex main && pdflatex main
\documentclass[11pt,a4paper]{article}

\usepackage[T1]{fontenc}
\usepackage[utf8]{inputenc}
\usepackage{lmodern}
\usepackage{microtype}
\usepackage[margin=25mm]{geometry}
\usepackage{amsmath,amssymb,bm}
\usepackage{graphicx}
\usepackage{subcaption}
\usepackage[font=small,labelfont=bf]{caption}
\usepackage{tikz}
\usetikzlibrary{arrows.meta}
\usepackage{booktabs}
\usepackage[hidelinks]{hyperref}
\usepackage[capitalise,nameinlink]{cleveref}
\usepackage{titlesec}
\usepackage[title]{appendix}

\numberwithin{equation}{section}
\setcounter{secnumdepth}{4}
\titleformat{\paragraph}{\normalfont\normalsize\itshape}{\theparagraph}{1em}{}
\titlespacing*{\paragraph}{0pt}{2ex plus 0.5ex minus 0.2ex}{1ex}

\newcommand{\Wsum}{\textstyle\sum_{j} W_j + W_b + W_0}   % common denominator
\newcommand{\pin}{p_{\mathrm{in}}}
\newcommand{\pout}{p_{\mathrm{out}}}
\newcommand{\patm}{p_{\mathrm{atm}}}
\newcommand{\plow}{p_{\mathrm{low}}}
\newcommand{\phigh}{p_{\mathrm{high}}}
\newcommand{\pth}{p_{\mathrm{th}}}
\newcommand{\R}{\mathbb{R}}

% Labelled figure: the image with TikZ labels placed in image coordinates (0..1 in x and y)
\newenvironment{labelledimage}[2][]{%
  \begin{tikzpicture}[#1]
  \node[anchor=south west, inner sep=0] (img) {\includegraphics[width=#2]{\labelledfile}};
  \begin{scope}[x={(img.south east)}, y={(img.north west)},
                lab/.style={font=\small, align=center, fill=white, fill opacity=0.85, text opacity=1,
                            rounded corners=1pt, inner sep=2pt},
                pin/.style={-{Stealth[length=2mm]}, thick}]}
  {\end{scope}\end{tikzpicture}}
\newcommand{\labelledfile}{}

\title{Developing a Feasible Neuroplastic Mechanical Neuron\\in the Domain of Soft Robotics}
\author{W.\,J. van Vuuren}  % TODO: co-authors and affiliation
\date{}

\begin{document}
\maketitle

\begin{abstract}
Onboard computation remains one of the open problems of soft robotics. This article develops a mechanically
feasible neuron that can serve as the building block of a neuroplastic mechanical neural network. Taking the
organisation of the biological neuron as inspiration, the proposed neuron encodes information in fluid pressure:
input pressures act on flexible membranes that bound a sealed activation chamber, so that the inputs are integrated
into a single field-based activation pressure without being consumed, and a pressure-controlled pneumatic valve
converts this activation pressure into a nonlinear output signal. The equilibrium of the membranes shows that the
activation pressure is a compliance-weighted combination of the input pressures, in which the mechanical weight of
each input is the ratio of its membrane area to its volumetric stiffness. Two mechanisms for changing these weights
after fabrication are proposed: limiting the deformation path of the membrane, and tuning the bulk modulus of an
intermediate fluid between two membranes. The sensitivity of the activation pressure to the chamber compliance, the
weight membranes, a bias membrane and the input pressures is derived. When the membrane stiffnesses are constant,
the neuron is exactly an artificial neuron, a weighted sum with a bias followed by a sigmoid, with weights that are
non-negative and, together with the bias and chamber weights, sum to one; the reachable hyperplanes therefore all
contain the equal-pressure diagonal, whose crossing point is set by the bias pressure. With deformation-dependent
stiffness, the neuron becomes a nonlinear, state-dependent computational element of which the artificial neuron is
a special case.
\end{abstract}

% =====================================================================================================================
\section{Introduction}
% =====================================================================================================================

\subsection{Background}

One of the remaining frontiers in soft robotics is the development of robust onboard computation and control. A
compelling long-term goal in this field is the creation of a fully mechanical robot capable of intelligent agency: a
system that can sense its environment, interact with it, and adapt autonomously based on those
interactions~\cite{rus2015soft}.

Nature has already achieved this through neural systems. Although biological organisms contain many forms of
computation, the brain remains the most sophisticated example of an integrated computational kernel. It is the
central neural network that underlies animal agency, and it is built from neural networks composed of simple
computational units known as neurons. In isolation, a single neuron is a modest computational element, but when
many neurons are interconnected, complex behaviour emerges. This broader principle, where intricate global
structure arises from the interaction of simple local rules, is widely observed in nature. Stephen Wolfram has
similarly speculated that physical reality itself may arise from repeated application of simple computational
rules, as developed in \emph{A New Kind of Science}~\cite{wolfram2002}.

The human brain contains more than 100~billion neurons, with each neuron potentially connecting to up to
10\,000 other neurons and passing signals through as many as 100~trillion synaptic
connections~\cite{zhang_neurons,stevens1979}. This vast network of interactions is one of the reasons the human
brain exhibits such remarkable computational capacity.

For this reason, there is a strong need to develop a mechanical neuron that can serve as a foundational building
block for mechanical neural networks (MNNs). Such a device would not merely mimic neural structure superficially; it
would need to support reprogrammable behaviour so that the resulting MNN could continue learning, adapting, and
reorganising in response to its environment. A true neuroplastic mechanical neuron would therefore provide the basis
for autonomous mechanical systems that do not simply react but evolve their responses over time.

\subsection{Objective and scope}

The objective of this work is to develop and characterise a mechanically feasible neuron that can serve as a
foundational building block for future neuroplastic mechanical neural networks. The proposed neuron is developed
from the physical organisation and functional roles of a biological neuron, with pressure-based mechanical systems
used to reproduce the processes of signal reception, integration, activation, and transmission. The intention is
therefore not to reproduce a biological neuron in its entirety, but to use its physical organisation as inspiration
for a mechanically feasible computational architecture.

The mechanical behaviour of the proposed neuron is subsequently formulated mathematically to establish how its
physical parameters determine its computational response. In particular, the analysis considers the mechanical
integration of multiple input pressures, the influence of membrane stiffness and geometry on the resulting
activation pressure, and the conversion of this activation into a nonlinear output. The ability to alter the
mechanical parameters responsible for input weighting is also investigated, providing a physical basis for
modifying the computational behaviour of the neuron after fabrication.

The resulting mechanical formulation is further related to the conventional artificial neuron. When the relevant
membrane stiffnesses are assumed to remain constant through deformation, the mechanical equilibrium can be
rearranged into a weighted-sum form with a bias term, and the nonlinear output relationship gives the corresponding
artificial-neuron formulation. This result provides a direct mathematical interpretation of the proposed mechanical
neuron in terms of an artificial neural-network neuron. However, constant stiffness represents a restricted case of
the more general mechanical system. Since membrane stiffness can vary with deformation, the effective mechanical
weights may themselves depend on the state of the neuron. The general formulation therefore describes a nonlinear,
state-dependent neural computation element, of which the conventional artificial neuron is a special case.

Finally, the mechanically adjustable parameters of the proposed architecture are considered in the context of
future neuroplasticity. This work does not implement an autonomous learning rule or demonstrate a complete
neuroplastic mechanical neural network. Rather, it establishes the mechanically reconfigurable neuron required as a
physical foundation for such a system. The development therefore progresses from biological inspiration to a
mechanically feasible architecture, from mechanical behaviour to its mathematical computational representation, and
finally to the reconfigurability required for future neuroplastic mechanical computation.

% =====================================================================================================================
\section{The neuron}
% =====================================================================================================================

\subsection{Neuron anatomy}

\begin{figure}[t]
  \centering
  \includegraphics[width=0.85\textwidth]{figures/neuron_anatomy.png}  % TODO: image credit if made with BioRender
  \caption{Anatomy of a biological neuron. The neuron consists of dendrites, which receive input signals from other
    neurons; the soma (cell body), where inputs are integrated; the axon, which transmits electrical signals to other
    neurons; and the synapse, where the axon communicates with the dendrites or soma of neighbouring neurons through
    chemical or electrical signals.}
  \label{fig:anatomy}
\end{figure}

This study focuses on the standard biological neuron and its major components, although alternative neuron
topologies exist. As illustrated in \cref{fig:anatomy}, a typical neuron consists of a soma, dendrites, a single
axon, and synapses~\cite{zhang_neurons}. While signal transmission in biological neurons is electrochemical in
nature, this work primarily focuses on the computational aspect of information processing (the kernel) rather than
the underlying mechanics.

Inputs to a neuron are received by the dendrites from other neurons. The dendrites process and integrate these
inputs through local computations~\cite{zhang_neurons,stevens1979}, which may involve complex interdependent
nonlinear interactions~\cite{london2005,beniaguev2021,tranvanminh2015,larkum2022}. These interactions are integrated
into a single value in the soma, which determines the neuron's response. When this integrated signal exceeds a
certain threshold, the neuron generates an action potential, a binary-like signal corresponding to a ``fire'' or
``no-fire'' state. This action potential then propagates along the axon to the synapses, which interface with the
dendrites of downstream neurons, allowing the signal to traverse the network~\cite{zhang_neurons,stevens1979}.

Although the magnitude of a neuron's action potential is effectively binary, the information it conveys is not
necessarily so. Neurons can encode information through temporal dynamics, most notably rate coding, in which signal
intensity is represented by the frequency of action potentials rather than their
amplitude~\cite{larkum2022,afshar2014,izhikevich2003}. This type of neural signalling is analogous to pulse-width or
pulse-frequency modulation: the information is not encoded in the spike magnitude but in the firing rate or spike
timing. As a result, rich, continuous information processing emerges from sequences of discrete, all-or-nothing
events.

From this perspective, the fundamental operations of a neuron can be conceptualised as a sequence of systematic
stages: input reception via synaptic--dendritic connections; dimensional collapse through inter-dendritic
computation into a unified neuronal activation within the soma; propagation of the resulting action potential along
the axon; and finally, transmission of this signal through synaptic connections to downstream neurons.

Accordingly, a neuron can be viewed as a system that receives multiple input signals, combines them into a single
activation through a parametric, potentially interdependent nonlinear relationship, and transmits the resulting
activation signal to the inputs of other neurons.

\subsection{Biological neural signal transmission}

Biological neurons transmit information through two complementary physical substrates, transport and field, that
operate in sequence across neuronal compartments. In dendrites, signalling is predominantly transport-based:
synaptic inputs open ion channels and generate local ionic currents that propagate through the cytoplasm toward the
soma, so information is conveyed by the movement of charge along dendritic cable-like structures~\cite{rall1994}.

As these currents reach the soma, the neuron shifts into a field-based regime of integration. In this regime,
changes in membrane potential regulate voltage-gated channels and conductances, enabling the cell to combine
multiple inputs through a local electrical state rather than through long-range ionic transport across the entire
compartment~\cite{koch2001,hodgkin1952}. This is the crucial computational step: the neuron does not simply pass the
input through unchanged but converts it into an internal state that determines whether and how an output will be
generated.

At chemical synapses, the process reverses. Depolarisation of the presynaptic membrane triggers vesicle fusion and
neurotransmitter release, which then diffuses across the cleft and produces transport-based ionic currents in the
postsynaptic cell~\cite{sudhof2013}. In this way, the neuron transforms an incoming transport-based signal into a
field-mediated internal response, and then back into a transport-based output.

This separation of transport-based sensing, field-mediated integration, and transport-based output is especially
important for mechanical neuron design because it enables non-destructive signal processing. The input is not
consumed by the act of computation; rather, it is sampled and transformed into a local state that can drive an
output response while preserving the integrity of the upstream signalling pathway. For artificial systems, this
distinction addresses a central engineering challenge: how to build a neuron that can compute on incoming
information without exhausting or destroying the signal that carries it.

\subsection{Neural plasticity}

Neural plasticity refers to the ability of a neuron to adapt by updating its internal parameters. In biological
neural networks, plasticity is a persistent and ongoing process: synaptic strengths are continuously modified as
organisms learn, adapt, and refine internal models of the world. This continual adaptation allows biological
systems to remain responsive to changing environments rather than converging to a fixed functional
state~\cite{demarin2014,mundkur2005}.

MNNs can be broadly classified according to their degree of neural plasticity. The first class comprises non-plastic
MNNs, in which the network parameters are determined prior to fabrication, often through analytical design,
numerical optimisation, or simulation, and are immutable once the system is manufactured. Such systems perform
computation through their mechanical structure but lack the ability to adapt or learn after deployment.

The second class consists of neuroplastic MNNs, in which network parameters can be modified after fabrication.
Within this class, further distinctions can be made based on the mechanism by which plasticity is achieved.
Mechanically neuroplastic MNNs implement parameter updates entirely through mechanical processes, enabling
adaptation without external domain intervention; the mechanical kernel itself is then solely responsible for the
parameter updating. In contrast, hybrid neuroplastic MNNs retain a mechanically embodied network structure but rely
on external systems to compute and apply parameter updates.

Representative examples of hybrid neuroplastic MNNs include architected materials that learn
behaviours~\cite{lee2022}, in which the mechanical network structure is adapted through electromechanical actuation,
and \emph{Beyond Silicon}~\cite{vanvuuren2025}, where parameter updates to the mechanical network are performed
manually.

% =====================================================================================================================
\section{Development of a mechanically feasible neuron}
% =====================================================================================================================

This section develops a conceptual mechanically feasible neuron by first defining the admissible design space and
the functional requirements imposed on the physical system. The design is initially formulated in an
implementation-agnostic manner, after which a topology inspired by the biological neuron is introduced. The
resulting model is then characterised mathematically to determine how its parameters influence signal transmission,
integration, thresholding, and plasticity.

\subsection{Design-space definition}

Before considering the physical structure of a mechanical neuron, it is necessary to define its computational
behaviour. Rather than prescribing a specific mechanical implementation, this definition establishes the minimum
computational requirements any proposed neuron must satisfy. Because signal processing in embodied systems relies
directly on physical mechanisms, computational and physical constraints are intrinsically interconnected.

The neuron is treated as a computational unit that receives multiple inputs, combines them according to a
parameter-dependent rule, and produces an output signal. A given neuron $i$ receives an input vector
\begin{equation}
  \vec{x}_i \in \R^{n_i},
  \label{eq:input-vector}
\end{equation}
where $n_i$ is the number of incoming signals. These inputs are received as transport-based signals, which are
transformed into a field-based scalar activation signal according to
\begin{equation}
  z_i = \sigma_i\bigl(f_i(\vec{x}_i, \bm{\theta}_i)\bigr) \in \R .
  \label{eq:neuron-map}
\end{equation}
In this expression, $f_i(\cdot, \bm{\theta}_i)$ represents the internal input-processing operation of the neuron.
It maps the $n_i$-dimensional transport-based input signal vector to a single field-based scalar quantity. The
parameter vector
\begin{equation}
  \bm{\theta}_i \in \Theta \subset \R^{p}
  \label{eq:parameters}
\end{equation}
contains the physical parameters that determine the behaviour of the neuron, which must be dynamically changeable
or reconfigurable to support plasticity. The function $\sigma_i(\cdot)$ represents the activation operation, which
maps the field-based scalar activation signal into a transport-based output signal via a nonlinear mapping to be
transmitted downstream.

\subsection{Biologically inspired mechanical neuron topology}

\subsubsection{Information encoding and signalling}

Computational information within the mechanical neuron is encoded in fluid pressure, with kernel operations
performed through the manipulation of pressure states. In an enclosed fluidic system, pressure is transmitted
through the substrate as a pressure field and does not require continuous fluid flow. This represents a field-based
signalling substrate, in which the pressure state is determined by the conditions throughout the enclosed fluid
volume. In a flow-based system, by contrast, the pressure depends on the dynamics of fluid moving through the
substrate. The signal is therefore transmitted through the transport of fluid along a channel and represents a
transport-based signalling substrate.

\subsubsection{Input-to-activation kernel}

\paragraph{Conceptual design of the flexible-membrane input interface}

\begin{figure}[t]
  \centering
  \renewcommand{\labelledfile}{figures/input_interface.png}
  \begin{labelledimage}{0.62\textwidth}
    \node[lab] at (0.215, 0.5) {Input fluid\\chamber};
    \node[lab] at (0.785, 0.5) {Input fluid\\chamber};
    \node[lab] at (0.5, 0.5) {Activation\\chamber};
    \node[lab] (m) at (0.5, 0.09) {Membrane};
    \draw[pin] (m.west) -- (0.338, 0.27);
    \draw[pin] (m.east) -- (0.649, 0.27);
  \end{labelledimage}
  \caption{Two-input configuration of the flexible-membrane input interface. Each input channel is separated from
    the central activation chamber by a flexible membrane.}
  \label{fig:input-interface}
\end{figure}

The input-to-activation stage of the neuron is implemented using flexible membranes that separate the input fluidic
channels from a central activation chamber. The activation chamber is filled with a confined fluid, while each
input channel conveys an externally generated pressure signal. Pressure applied within an input channel causes its
corresponding membrane to deform, thereby transmitting a force to the fluid contained in the central chamber. The
contributions from the individual input pressures are consequently combined through the mechanical interaction of
the membranes and the activation fluid. The resulting pressure within the central chamber constitutes the neuron's
scalar activation signal.

\Cref{fig:input-interface} illustrates a two-input configuration of this mechanism. Changes in the input pressures
alter the forces acting on the respective membranes, which in turn modify the pressure of the fluid contained within
the activation chamber. The chamber pressure therefore provides a single internal state that reflects the combined
influence of the input signals.

This configuration also provides a natural means of achieving non-destructive signal processing. The input signals
are supplied by pressure sources operating at a prescribed steady-state pressure with continuous flow through the
input channels. In contrast, the fluid within the activation chamber is enclosed and does not flow into or out of
the chamber during operation. The input pressures therefore influence the activation state through
membrane-mediated force transmission, rather than through the transfer of fluid into the activation chamber. The
input signals can consequently be processed without being consumed or directly altered by the activation
mechanism.

The proposed arrangement thus implements the first stage of the mechanical neuron's computational operation:
multiple transport-based input signals are converted into a single field-based activation signal.

\paragraph{Mathematical model of the flexible-membrane input interface}

The following characterises the mechanical behaviour of the membranes and activation chamber and establishes how
the input pressures determine the resulting chamber pressure.

At equilibrium, each flexible membrane deforms until the elastic restoring force generated by its deformation
balances the net pressure force acting across it. The pressure difference across a membrane is determined by the
pressure in the corresponding input channel and the pressure in the central activation chamber.

For membrane $j$, let $p_j$ denote the pressure in input channel $j$, $p_a$ the pressure in the activation chamber,
and $\patm$ the atmospheric pressure. If $A_j$ is the effective area of the membrane and $F_j^{e}$ is its elastic
restoring force, the equilibrium condition is
\begin{equation}
  F_j^{e} = (p_j - p_a)\,A_j .
  \label{eq:membrane-equilibrium}
\end{equation}
This expression assumes a sign convention in which positive membrane deformation is produced by an input pressure
greater than the activation-chamber pressure. More generally, the membrane equilibrium may be written as
\begin{equation}
  F_j^{e}(\delta_j) - (p_j - p_a)\,A_j = 0,
  \label{eq:membrane-equilibrium-general}
\end{equation}
where $\delta_j$ is the equilibrium deformation of membrane $j$. The elastic force is therefore a function of
deformation, and the deformation is determined implicitly by the pressure difference across the membrane.

To relate membrane deformation to its restoring force, the deformation is represented in terms of the volume
displaced by the membrane rather than its local displacement field. Let $\Delta V_j$ denote the volumetric
deformation of membrane $j$, defined as the change in activation-chamber volume produced by its deformation. A
volumetric stiffness $K_j$ can then be introduced to describe the resistance of the membrane to this deformation.
The elastic restoring force is expressed as
\begin{equation}
  F_j^{e} = K_j(\cdot)\,\Delta V_j .
  \label{eq:volumetric-stiffness}
\end{equation}
This represents a generalised stiffness relation, in which the stiffness may vary with the current volumetric
deformation. Substituting it into the membrane equilibrium gives
\begin{subequations}
\begin{align}
  K_j(\cdot)\,\Delta V_j &= (p_j - p_a)\,A_j, \label{eq:dv-a}\\
  \Delta V_j &= \frac{(p_j - p_a)\,A_j}{K_j(\cdot)} . \label{eq:dv-b}
\end{align}
\end{subequations}
Under the assumption that the activation fluid is incompressible and the chamber is sealed, the total change in
fluid volume must be zero. Including the $m$ input membranes, the condition is
\begin{equation}
  \sum_{j=1}^{m} \Delta V_j = 0 .
  \label{eq:volume-conservation}
\end{equation}
Substituting the membrane deformation relations gives
\begin{equation}
  \sum_{j=1}^{m} \frac{(p_j - p_a)\,A_j}{K_j(\cdot)} = 0 .
  \label{eq:volume-balance}
\end{equation}
For constant effective stiffnesses, solving for the activation-chamber pressure yields
\begin{equation}
  p_a = \frac{\displaystyle\sum_{j=1}^{m} \frac{A_j}{K_j(\cdot)}\,p_j}{\displaystyle\sum_{j=1}^{m} \frac{A_j}{K_j(\cdot)}} .
  \label{eq:pa-compliance}
\end{equation}
Thus, the equilibrium activation pressure is a stiffness- and area-weighted combination of the input pressures. The
weighting of each input is determined by its membrane's effective compliance, $A_j/K_j(\cdot)$. Since this quantity
determines the relative contribution of each input pressure to the activation pressure, it is analogous to a weight
in an artificial neural network. The effective mechanical weight of membrane $j$ is therefore defined as
\begin{equation}
  W_j(\cdot) = \frac{A_j}{K_j(\cdot)} .
  \label{eq:weight}
\end{equation}
Using this definition, the activation pressure becomes
\begin{equation}
  p_a = \frac{\sum_{j=1}^{m} W_j(\cdot)\,p_j}{\sum_{j=1}^{m} W_j(\cdot)} .
  \label{eq:pa-weights}
\end{equation}
Since $W_j(\cdot) > 0$, the activation pressure is a convex combination of the input pressures.

\subsubsection{Mechanically reconfigurable weights}

\paragraph{Weight-modification principle}

The mechanical weight associated with each input must be physically adjustable for the neuron to exhibit
neuroplastic behaviour. Since the weight is defined by the ratio between the effective membrane area and stiffness,
it can be modified by changing either of these properties. The proposed approach focuses on modifying the effective
stiffness through controlled deformation of the input membrane.

The stiffness of membrane $j$ is not treated as a fixed quantity. Instead, it is considered a function of the
membrane's equilibrium deformation. The deformation itself depends on the pressure difference across the membrane,
although it may also be influenced by other physical, geometric, material, and boundary-condition variables. These
dependencies are represented as
\begin{equation}
  K_j = K_j(\delta_j, \cdot)
  \label{eq:k-of-delta}
\end{equation}
and
\begin{equation}
  \delta_j = \delta_j(p_j - p_a, \cdot),
  \label{eq:delta-of-dp}
\end{equation}
where the notation $(\cdot)$ represents additional variables that may influence the respective relationship.
Substituting these dependencies into the definition of the mechanical weight gives
\begin{equation}
  W_j = \frac{A_j}{K_j\bigl(\delta_j(p_j - p_a, \cdot), \cdot\bigr)} .
  \label{eq:weight-composed}
\end{equation}
The weight may therefore be represented more generally as
\begin{equation}
  W_j = W_j(p_j - p_a, \cdot) .
  \label{eq:weight-state}
\end{equation}
This relationship demonstrates that the mechanical weight is coupled to the equilibrium state of the membrane. By
controlling the deformation of membrane $j$, its effective stiffness can be altered, and its corresponding weight can
consequently be updated. The pressure difference across the membrane provides one means of influencing this
deformation, but it is not the sole determinant of the weight. Rather, the weight depends on the complete mechanical
state of the membrane and the other variables governing its deformation and stiffness.

\paragraph{Proposed mechanical realisation of the weight update}

\subparagraph{Deformation limiting.}

\begin{figure}[t]
  \centering
  \begin{subfigure}[b]{0.36\textwidth}
    \centering
    \includegraphics[width=\textwidth]{figures/limit_outward.png}
    \caption{Outward deformation limiting}
    \label{fig:limit-outward}
  \end{subfigure}\hspace{0.06\textwidth}%
  \begin{subfigure}[b]{0.36\textwidth}
    \centering
    \includegraphics[width=\textwidth]{figures/limit_inward.png}
    \caption{Inward deformation limiting}
    \label{fig:limit-inward}
  \end{subfigure}
  \caption{Methods for altering membrane deformation. (a)~Outward membrane deformation restriction, limiting the
    membrane's deformation away from the activation chamber. (b)~Inward membrane deformation restriction, limiting
    the membrane's displacement into the activation chamber.}
  \label{fig:deformation-limiting}
\end{figure}

Deformation-limited weight tuning operates by restricting the membrane's available deformation path. Since the
membrane can deform in either direction, the limiting obstacle may be positioned either on the activation-chamber
side, producing inward deformation limiting, or on the input-channel side, producing outward deformation limiting,
as shown in \cref{fig:deformation-limiting}.

By controlling the shape of the limiting obstacle, the allowable equilibrium deformation of the membrane can be
controlled. Deformation can be limited independently in either direction, such that
\begin{equation}
  \delta_j =
  \begin{cases}
    \delta_j^{\mathrm{in}},  & p_a < p_j,\\[2pt]
    \delta_j^{\mathrm{out}}, & p_a > p_j .
  \end{cases}
  \label{eq:deformation-limits}
\end{equation}
Here, $\delta_j^{\mathrm{in}}$ and $\delta_j^{\mathrm{out}}$ denote the inward and outward deformation limits,
respectively. Since the effective stiffness is a function of the membrane deformation, $K_j = K_j(\delta_j, \cdot)$,
the mechanical weight is consequently also affected:
\begin{equation}
  W_j = \frac{A_j}{K_j(\delta_j, \cdot)} .
  \label{eq:weight-limited}
\end{equation}
Deformation limiting therefore provides a mechanical means of updating the weight. Changing the deformation path
changes the membrane's equilibrium deformation and effective stiffness, which consequently changes the contribution
of the corresponding input to the activation pressure.

Inward deformation limiting is expected to introduce an additional complication because the limiting obstacle would
be located inside the activation chamber. Changing its volume changes the volume available to the confined
activation fluid, coupling the weight-update mechanism to the chamber's pressure. This complication is avoided when
the deformation limit is implemented on the input-channel side.

\subparagraph{Intermediate-fluid bulk-modulus tuning.}

\begin{figure}[t]
  \centering
  \renewcommand{\labelledfile}{figures/bulk_modulus_weight.png}
  \begin{labelledimage}{0.78\textwidth}
    \node[lab, fill opacity=0] (i1) at (0.183, 1.10) {Input};
    \node[lab, fill opacity=0] (ac) at (0.5, 1.10) {Activation chamber};
    \node[lab, fill opacity=0] (i2) at (0.817, 1.10) {Input};
    \node[lab, fill opacity=0] (w1) at (0.25, -0.12) {Weight chamber};
    \node[lab, fill opacity=0] (w2) at (0.75, -0.12) {Weight chamber};
    \draw[pin] (i1.south) -- (0.183, 0.86);
    \draw[pin] (ac.south) -- (0.5, 0.86);
    \draw[pin] (i2.south) -- (0.817, 0.86);
    \draw[pin] (w1.north) -- (0.30, 0.20);
    \draw[pin] (w2.north) -- (0.70, 0.20);
  \end{labelledimage}
  \caption{Conceptual schematic of an intermediate-fluid bulk-modulus-tuned input-weight kernel: flexible
    membranes separated by an intermediate (weight) chamber whose tuneable compressibility controls the transmission
    of deformation and force from the input-pressure chamber to the activation-pressure chamber.}
  \label{fig:bulk-modulus}
\end{figure}

Force-transfer limiting is implemented by dividing the original membrane into two flexible membranes with a
confined fluid layer between them (\cref{fig:bulk-modulus}). The input pressure acts on the outer membrane, while the
activation-chamber pressure acts on the inner membrane. The pressure of the intermediate fluid, referred to as the
weight pressure, transmits deformation between the two membranes.

This configuration introduces three pressure variables: the input pressure $p_i$, the activation-chamber pressure
$p_a$, and the weight pressure $p_w$. Let $\Delta V_1$ denote the deformation of membrane~1 into the intermediate
chamber, with $K_1$ representing its volumetric stiffness. Similarly, let $\Delta V_2$ denote the deformation of
membrane~2 into the activation chamber, with $K_2$ representing its volumetric stiffness.

It is important to note that $\Delta V_j$ is defined as the change in the volume of the activation chamber.
Consequently, $K_j$ represents the volumetric stiffness associated with deformation into the activation chamber due
to the difference between the input pressure $p_i$ and the activation pressure $p_a$. The detailed derivation of
$F_j$, $K_j$, and their dependence on the intermediate-fluid properties is provided in Appendix~\ref{app:bulk}.

The intermediate fluid provides the tuneable mechanical element. Its bulk modulus $B$ determines its resistance to
volumetric compression and therefore controls how effectively deformation is transferred from membrane~1 to
membrane~2. A low bulk modulus allows the intermediate fluid to undergo a comparatively large volume change,
reducing force transfer to membrane~2. In the limiting case $B = 0$, the intermediate fluid provides no resistance to
compression, and the effective stiffness of the combined mechanism approaches
\begin{equation}
  K_j(B = 0) \rightarrow \infty .
  \label{eq:k-b0}
\end{equation}
In contrast, as the bulk modulus increases, the intermediate fluid becomes less compressible and transfers an
increasing proportion of the deformation between the two membranes. In the limiting case $B \rightarrow \infty$, the
intermediate fluid is incompressible, and the deformation of the two membranes becomes fully coupled. The effective
stiffness then approaches the sum of the individual membrane stiffnesses:
\begin{equation}
  K_j(B \rightarrow \infty) = K_1 + K_2 .
  \label{eq:k-binf}
\end{equation}
The effective stiffness can therefore be represented as a function of the bulk modulus,
\begin{equation}
  K_j = K_j(B),
  \label{eq:k-of-b}
\end{equation}
with the limiting behaviour
\begin{equation}
  K_j(0) = \infty, \qquad K_j(\infty) = K_1 + K_2 .
  \label{eq:k-limits}
\end{equation}
Since the mechanical weight is inversely related to the effective stiffness, it follows that
\begin{equation}
  W_j(B) = \frac{A_j}{K_j(B)} .
  \label{eq:weight-of-b}
\end{equation}
The bulk modulus of the intermediate fluid therefore provides a physical means of tuning the mechanical weight. By
changing $B$, the force-transfer capability of the two-membrane system is altered, which changes the effective
stiffness $K_j$ and consequently the influence of the corresponding input on the activation pressure. Bulk modulus
is conventionally defined as the pressure change associated with volumetric strain, supporting this interpretation of
$B$ as the fluid's resistance to compression.

\subsubsection{Activation-to-output kernel}

The activation pressure $p_a$ is a field-based internal state. To communicate this state to another mechanical
neuron, the activation pressure must act on an output interface that generates a pressure-driven flow or pressure
pulse in an external output channel.

The proposed design uses a pressure-controlled pneumatic switching interface to perform this conversion. The
interface consists of four fundamental components: a pressure source, an atmospheric vent, a flow resistor, and a
mechanically actuated valve. The pressure source establishes the high-pressure state of the output signal, while the
atmospheric vent establishes the low-pressure state. The flow resistor restricts fluid flow. The mechanically
actuated valve is controlled by the activation pressure $p_a$, opening or closing according to whether the
activation state is low or high.

\paragraph{Inverted output configuration}

\begin{figure}[t]
  \centering
  \begin{subfigure}[b]{0.48\textwidth}
    \centering
    \includegraphics[width=\textwidth]{figures/output_inverted.png}
    \caption{Inverted configuration}
    \label{fig:output-inverted}
  \end{subfigure}\hfill
  \begin{subfigure}[b]{0.48\textwidth}
    \centering
    \includegraphics[width=\textwidth]{figures/output_noninverted.png}
    \caption{Non-inverted configuration}
    \label{fig:output-noninverted}
  \end{subfigure}
  \caption{Activation-to-output configurations with a normally open, pressure-controlled mechanical valve driven by
    $p_a$. (a)~Inverted: the valve sits between the pressure source and the output connection, and the output is
    connected through a flow resistor to atmosphere, so that a low activation pressure maintains a high output
    pressure and a high activation pressure produces a low output pressure. (b)~Non-inverted: the flow resistor sits
    between the pressure source and the output connection, and the valve between the output and the atmospheric
    vent, so that a low activation pressure produces a low output pressure and a high activation pressure produces a
    high output pressure.}
  \label{fig:output}
\end{figure}

The inverted output configuration is illustrated in \cref{fig:output-inverted}. In this configuration the
pressure-controlled mechanical valve is on the pressure-source side, and the flow resistor is on the vent side. The
output connection to the downstream neuron inputs is located between the flow resistor and the valve.

Assume the pressure valve is normally open. When the activation pressure $p_a$ is low, the valve remains open and
connects the output to the pressure source. Fluid continuously escapes through the flow resistor to atmosphere, but
this loss is replenished by the pressure source. The output, which is connected to neuron inputs that do not draw
fluid flow, therefore remains at a high pressure. When $p_a$ is high, the valve closes and disconnects the output
from the pressure source. The output remains connected to atmosphere through the flow resistor, allowing fluid to
escape without being replenished. The output pressure therefore decreases to the low-pressure state. Consequently,
the configuration produces an inverted output relationship:
\begin{equation*}
  p_a \text{ low} \Rightarrow \pout \text{ high}, \qquad
  p_a \text{ high} \Rightarrow \pout \text{ low}.
\end{equation*}

\paragraph{Non-inverted output configuration}

The non-inverted output configuration is illustrated in \cref{fig:output-noninverted}. In this configuration, the
flow resistor is positioned on the pressure-source side, while the pressure-controlled mechanical valve is positioned
on the vent side. The output connection to the downstream neuron inputs is located between the flow resistor and
the valve.

Assuming again that the pressure valve is normally open, a low activation pressure $p_a$ keeps the valve open and
connects the output to atmosphere. The output pressure can therefore escape through the valve, while the pressure
source replenishes the output through the flow resistor. Because the replenishment flow is restricted, the output
pressure decreases to the low-pressure state. When $p_a$ is high, the valve closes and disconnects the output from
atmosphere. The output remains connected to the pressure source through the flow resistor, allowing its pressure to
increase towards the source pressure. Once the valve is closed, the output pressure cannot escape through the vent
and therefore remains at the high-pressure state. Consequently, the configuration produces a non-inverted output
relationship:
\begin{equation*}
  p_a \text{ low} \Rightarrow \pout \text{ low}, \qquad
  p_a \text{ high} \Rightarrow \pout \text{ high}.
\end{equation*}

% =====================================================================================================================
\section{Mathematical model of the mechanical neuron}
% =====================================================================================================================

This section first characterises the mechanical behaviour of the proposed neuron by examining the effects of the
activation-chamber stiffness, the membrane stiffness, the bias membrane term and the input pressures. This
characterisation establishes how the physical properties and initial operating conditions of the neuron influence
its response. Building on the mechanical characterisation, the analysis then establishes the computational
expressiveness of the proposed neuron.

\subsection{Mathematical model}

The parameters $\plow$ and $\phigh$ define the lower and upper pressure bounds of the system, respectively. These
bounds represent the minimum and maximum pressures considered.

Using \cref{eq:pa-weights}, and singling out the first input, the activation pressure of the neuron is written as
\begin{equation}
  p_a = \frac{W_1 p_1 + \sum_{j\neq 1} W_j p_j + W_b p_b + W_0 \patm}{W_1 + \sum_{j\neq 1} W_j + W_b + W_0} .
  \label{eq:pa-full}
\end{equation}
Here, $p_j$ is the $j$-th input pressure, while $W_j$ is the corresponding mechanical weight. The term $W_b p_b$
acts as a bias, supplied by a bias membrane with weight $W_b$ and bias pressure $p_b$. The term $W_0$ represents the
compliance of the activation chamber with respect to the atmosphere, and $\patm$ is the atmospheric pressure acting
on the chamber housing.

The activation-to-output kernel is represented by a sigmoid function:
\begin{equation}
  \pout(p_a) = \plow + \frac{\phigh - \plow}{1 + \exp\bigl[-k\,(p_a - \pth)\bigr]} .
  \label{eq:sigmoid}
\end{equation}
Here, $p_a$ is the activation pressure, $\pth$ is the switching threshold, $\plow$ and $\phigh$ are the limiting
output pressures, and $k$ controls the steepness and direction of the transition. A positive $k$ represents the
non-inverted configuration, in which the output pressure increases with $p_a$, whereas a negative $k$ represents the
inverted configuration, in which the output pressure decreases as $p_a$ increases.

For simplicity, unless stated otherwise, all subsequent analyses assume that the membrane areas and stiffnesses
remain constant. The membranes are therefore modelled as linear elastic elements, such that their restoring forces
vary linearly with deformation.

\subsection{Mechanical characterisation}

For the mechanical characterisation, atmospheric pressure is used as the reference pressure, $\patm = 0$, and the
pressure range considered in the model is bounded by
\begin{equation*}
  \plow = -100~\mathrm{kPa}, \qquad \phigh = 100~\mathrm{kPa}.
\end{equation*}

\subsubsection{Activation-chamber stiffness}

The activation-chamber compliance is represented by $W_0$. With $\patm = 0$, \cref{eq:pa-full} becomes
\begin{equation}
  p_a = \frac{\sum_j W_j p_j + W_b p_b}{\sum_j W_j + W_b + W_0} .
  \label{eq:pa-atm0}
\end{equation}
Defining the weighted input pressure as
\begin{equation}
  \pin = \sum_j W_j p_j + W_b p_b,
  \label{eq:pin}
\end{equation}
the activation pressure becomes
\begin{equation}
  p_a = \frac{1}{\sum_j W_j + W_b + W_0}\,\pin .
  \label{eq:pa-pin}
\end{equation}
Let
\begin{equation}
  \alpha = \frac{1}{\sum_j W_j + W_b + W_0} .
  \label{eq:alpha}
\end{equation}
Then $p_a = \alpha\,\pin$. Substituting this relationship into \cref{eq:sigmoid} gives
\begin{equation}
  \pout(\pin) = \plow + \frac{\phigh - \plow}{1 + \exp\bigl[-k\,(\alpha\,\pin - \pth)\bigr]} .
  \label{eq:sigmoid-pin}
\end{equation}
Thus, the activation-chamber compliance is incorporated into the activation function through the scaling factor
$\alpha$. When the function is expressed in terms of the weighted input pressure, this factor modifies the
effective gradient of the sigmoid:
\begin{equation}
  k_{\mathrm{eff}} = k\,\alpha .
  \label{eq:keff}
\end{equation}
Therefore, increasing the chamber compliance reduces the effective gradient of the activation-to-output function.
The complete algebraic transformation is provided in Appendix~\ref{app:compliance}.

\subsubsection{Weight membrane}

\paragraph{Sensitivity}

To examine the influence of the first weight membrane, $W_1$ is treated as the variable parameter while the
remaining pressures and weights are held constant. The activation pressure is
\begin{equation}
  p_a = \frac{W_1 p_1 + \sum_{j\neq 1} W_j p_j + W_b p_b + W_0 \patm}{W_1 + \sum_{j\neq 1} W_j + W_b + W_0} .
  \label{eq:pa-w1}
\end{equation}
The sensitivity of the activation pressure to the mechanical weight $W_1$ is (Appendix~\ref{app:weight-sensitivity})
\begin{equation}
  \frac{\partial p_a}{\partial W_1} = \frac{p_1 - p_a}{W_1 + \sum_{j\neq 1} W_j + W_b + W_0} .
  \label{eq:dpa-dw1}
\end{equation}
Thus, increasing $W_1$ shifts the activation pressure towards $p_1$. The corresponding sensitivity to the membrane
stiffness $K_1$, using $W_1 = A_1/K_1$, is
\begin{equation}
  \frac{\partial p_a}{\partial K_1} =
    -\frac{p_1 - p_a}{K_1 + \dfrac{K_1^2}{A_1}\Bigl(\sum_{j\neq 1} W_j + W_b + W_0\Bigr)} .
  \label{eq:dpa-dk1}
\end{equation}
This expression shows how changes in the membrane stiffness affect the activation pressure. Because the mechanical
contribution associated with the membrane is proportional to $A_1/K_1$, decreasing $K_1$ increases the relative
contribution of $p_1$ to $p_a$.

\paragraph{Range}

Since $W_1 = A_1/K_1$, the limiting cases are
\begin{equation}
  \lim_{K_1 \rightarrow 0} p_a = p_1,
  \label{eq:lim-k0}
\end{equation}
because $W_1 \rightarrow \infty$ and the contribution of $p_1$ dominates the weighted average. Conversely,
\begin{equation}
  \lim_{K_1 \rightarrow \infty} p_a =
    \frac{\sum_{j\neq 1} W_j p_j + W_b p_b + W_0 \patm}{\sum_{j\neq 1} W_j + W_b + W_0},
  \label{eq:lim-kinf}
\end{equation}
because $W_1 \rightarrow 0$. Thus, a membrane with zero stiffness makes $p_a$ approach $p_1$, whereas an infinitely
stiff membrane makes its contribution negligible.

\subsubsection{Bias membrane term}

Starting from \cref{eq:pa-atm0}, and holding all other parameters constant, the quotient rule gives
\begin{equation}
  \frac{\partial p_a}{\partial W_b} = \frac{p_b - p_a}{\sum_j W_j + W_b + W_0} = \alpha\,(p_b - p_a),
  \label{eq:dpa-dwb}
\end{equation}
where $\alpha$ is defined in \cref{eq:alpha}. Similarly,
\begin{equation}
  \frac{\partial p_a}{\partial p_b} = \frac{W_b}{\sum_j W_j + W_b + W_0} = \alpha\,W_b .
  \label{eq:dpa-dpb}
\end{equation}
When $W_b$ is held constant, $\partial p_a/\partial p_b = \alpha W_b$ remains constant for fixed system parameters.
Therefore, $p_b$ acts as a predictable and tuneable bias input. In contrast, holding $p_b$ constant does not make
$\partial p_a/\partial W_b = \alpha(p_b - p_a)$ constant, because this sensitivity depends on the changing value of
$p_a$. Consequently, $W_b$ does not act as a static bias.

\subsubsection{Input pressure and ANN form}

\paragraph{Input sensitivity}

The input pressures $p_j$ are the externally applied pressure signals that drive the mechanical neuron. Each input
acts on a membrane with an associated mechanical weight $W_j$, and the resulting pressure contributions are combined
in the activation chamber. The purpose of this section is to determine how each input pressure influences the
activation pressure and to establish the relationship between the mechanical model and the mathematical form of an
artificial neural network (ANN).

The sensitivity of the activation pressure to an individual input pressure $p_i$ is (Appendix~\ref{app:input-sensitivity})
\begin{equation}
  \frac{\partial p_a}{\partial p_i} = \frac{W_i}{\sum_j W_j + W_b + W_0} .
  \label{eq:dpa-dpi}
\end{equation}
Using the definition of $\alpha$,
\begin{equation}
  \frac{\partial p_a}{\partial p_i} = \alpha\,W_i .
  \label{eq:dpa-dpi-alpha}
\end{equation}
This result shows that the influence of each input pressure is determined by its corresponding mechanical weight
relative to the total effective weight of the system. A larger $W_i$ produces a larger change in $p_a$ for a given
change in $p_i$, while a smaller $W_i$ produces a weaker response.

\paragraph{Special case: ANN form}

Assuming that $W$ is constant for all membranes, the activation pressure can be written in the same form as an
artificial neuron,
\begin{equation}
  p_a = \sum_j \left(\frac{\partial p_a}{\partial p_j}\right) p_j + b,
  \label{eq:ann-sensitivity-form}
\end{equation}
where $b$ represents the artificial bias contribution. From the derivation in Appendix~\ref{app:ann}, the effective
artificial input weights and bias are
\begin{equation}
  w_j = \frac{\partial p_a}{\partial p_j} = \alpha\,W_j
  \label{eq:ann-weight}
\end{equation}
and
\begin{equation}
  b = \alpha\,W_b\,p_b = w_b\,p_b .
  \label{eq:ann-bias}
\end{equation}
Thus,
\begin{equation}
  p_a = \sum_j w_j\,p_j + b .
  \label{eq:ann-preactivation}
\end{equation}
This is the standard pre-activation form of an artificial neuron. In this representation, the pressure $p_j$
corresponds to the neural-network input, the normalised mechanical weight $\alpha W_j$ corresponds to the ANN weight,
and the bias-membrane contribution $\alpha W_b p_b$ corresponds to the ANN bias.

The activation-to-output relationship is then applied to the activation pressure,
\begin{equation}
  \pout = F(p_a),
  \label{eq:af}
\end{equation}
where $F$ is the sigmoid pressure-response function defined in \cref{eq:sigmoid}. Substituting the equivalent ANN
expression for $p_a$ gives
\begin{equation}
  \pout = F\Bigl(\sum_j w_j\,p_j + b\Bigr) .
  \label{eq:ann-full}
\end{equation}
This establishes that the complete input-to-output kernel of the proposed mechanical neuron is computationally
equivalent to that of an artificial neuron when the membrane areas and effective stiffnesses remain constant
throughout deformation. In this case, the sensitivity of the activation pressure to each input pressure is constant,
allowing the activation pressure to be expressed as a weighted sum of the inputs and a bias term.

\paragraph{General form}

This, however, represents a special case of the more general mechanical formulation. In general, the membrane
stiffness varies with deformation, and the deformation depends on both the activation pressure and the corresponding
input pressure. Consequently, the sensitivity of the activation pressure to an input is generally state-dependent,
\begin{equation}
  \frac{\partial p_a}{\partial p_i} =
    \frac{w_i + \sum_j p_j \dfrac{\partial w_j}{\partial p_i} + \dfrac{\partial b}{\partial p_i}}
         {1 - \sum_j p_j \dfrac{\partial w_j}{\partial p_a} - \dfrac{\partial b}{\partial p_a}},
  \label{eq:general-sensitivity}
\end{equation}
and the effective mechanical weight can therefore be expressed as a function of the instantaneous state and input.
The activation pressure consequently takes the more general form
\begin{equation}
  p_a = \sum_j w_j(p_a, \vec{p}, p_b)\,p_j + b(p_a, \vec{p}, p_b), \qquad
  \vec{p} = [p_1, p_2, \ldots, p_N]^{\mathsf{T}} .
  \label{eq:general-form}
\end{equation}
The neuron is therefore not, in general, equivalent to a conventional fixed-weight artificial neuron. Instead, it
constitutes a nonlinear, state-dependent neural computation element.

\subsection{Computational expressiveness}

The computational expressiveness of the proposed neuron is determined by the space of parameters that can be
physically represented by its mechanical architecture. In the constant-stiffness case, where the neuron can be
expressed in artificial-neuron form, the input sensitivity is given by \cref{eq:dpa-dpi}, while the corresponding
bias is
\begin{equation*}
  b = \frac{W_b\,p_b}{\sum_j W_j + W_b + W_0} .
\end{equation*}
Consequently, the effective ANN parameters satisfy $\sum_j w_j + w_b + w_0 = 1$. Furthermore, since the underlying
membrane areas and stiffnesses cannot be negative, the corresponding mechanical weights are constrained to be
positive. The resulting neuron therefore occupies only a restricted subset of the parameter space of a conventional
artificial neuron.

The expressiveness is first examined within this artificial-neuron representation, corresponding to the special
case of constant stiffness. The analysis then considers the general case, in which membrane stiffness varies with
deformation and the effective weights become functions of the activation and input pressures. There, the neuron is
treated as a nonlinear, state-dependent computational element, and its expressiveness is considered within this
more general parameter space.

\subsubsection{Constant-stiffness computational expressiveness}

For the constant-stiffness case, the proposed neuron can be represented by the artificial-neuron formulation
\begin{equation}
  p_a = \sum_j w_j\,p_j + w_b\,p_b + w_0\,p_0 .
  \label{eq:ann-with-reference}
\end{equation}
The parameters $w_j$, $w_b$, and $p_b$ are mechanically tuneable, whereas $w_0$ and $p_0$ are static parameters
representing the compliance of the activation chamber to the atmospheric reference. Consequently, an $N$-input
mechanical neuron has $N+2$ tuneable parameters: $N$ input weights, the mechanical bias weight $w_b$, and the
tuneable bias pressure $p_b$. In contrast, a conventional $N$-input artificial neuron has $N+1$ parameters,
consisting of $N$ weights and a single independent bias term. The additional mechanical parameter arises because the
bias contribution is physically represented by the product $w_b p_b$, rather than by a single independently
specified parameter.

The pre-activation relation defines a hyperplane in the input--output space, with the weights determining its
orientation. Defining $w_j = \tan\theta_j$, each weight determines the inclination $\theta_j$ of the hyperplane
relative to the corresponding input axis. Since the mechanical parameters from which the weights are formed are
non-negative, $w_j \geq 0$, and therefore $\theta_j \geq 0$. Consequently,
\begin{equation}
  \frac{\partial p_a}{\partial p_j} = w_j \geq 0,
  \label{eq:monotone}
\end{equation}
meaning that increasing any input pressure cannot decrease the pre-activation pressure. The mechanical neuron is
therefore monotonically non-decreasing with respect to each input, and strictly increasing where $w_j > 0$. The
additional constraint
\begin{equation}
  \sum_j w_j + w_b + w_0 = 1
  \label{eq:normalisation}
\end{equation}
further restricts the representable parameter space by removing one degree of freedom from the weights. This
constraint also imposes a geometric property on the resulting hyperplane: it must contain the equal-pressure
diagonal, defined by
\begin{equation}
  p_1 = p_2 = \cdots = p_N = p_b = p_0 = p_a .
  \label{eq:diagonal}
\end{equation}
As shown in Appendix~\ref{app:hyperplane}, fixing $p_b$ reduces this diagonal to a specific point through which the
hyperplane must pass. Since $p_b$ is itself tuneable, this point can be moved along the equal-pressure diagonal,
providing an additional degree of freedom in the position of the hyperplane. Thus, although the constant-stiffness
mechanical neuron has $N+2$ tuneable mechanical parameters, the normalisation constraint leaves $N+1$ independent
degrees of freedom, matching the number of independent parameters of a conventional $N$-input artificial neuron.
However, the resulting parameter space remains restricted by the positivity and normalisation constraints and
therefore does not span the full parameter space of a conventional artificial neuron.

\subsubsection{General computational expressiveness}

For the general case, the normalisation constraint in \cref{eq:normalisation} remains present. Therefore, the
derivation presented in Appendix~\ref{app:hyperplane} remains applicable, and the equal-pressure diagonal continues to
satisfy the general formulation. For a fixed $p_b$, the computational surface must consequently pass through the
point on the diagonal corresponding to $p_b$,
\begin{equation*}
  (p_1, \ldots, p_N, p_a) = (p_b, \ldots, p_b, p_b) .
\end{equation*}
Thus, $p_b$ continues to determine the point at which the computational surface intersects the equal-pressure
diagonal, while the state-dependent weights determine the local shape and orientation of the surface around this
point.

The positivity constraint on the mechanical parameters also remains. Although the membrane stiffness is no longer
constant, its value remains physically positive,
\begin{equation*}
  K_i(\delta_i) > 0,
\end{equation*}
for all admissible deformations $\delta_i$. Consequently, the effective mechanical weights remain within the
positive domain. However, unlike the constant-stiffness case, $K_i(\delta_i)$ is generally neither constant nor
necessarily monotonic with deformation. The effective weights therefore vary with the mechanical state, and the
resulting input-to-output mapping is not necessarily monotonically increasing with each input. The general neuron
consequently remains constrained to the physically positive mechanical parameter domain while allowing a nonlinear,
state-dependent computational surface within that domain.

% =====================================================================================================================
\section{Conclusion}
% =====================================================================================================================

This article established the theoretical, structural, and mathematical framework for a mechanically feasible
neuroplastic neuron designed for soft robotic applications. By translating biological neural processes into fluidic
mechanisms, the proposed design demonstrates how pressure-driven inputs can be integrated through membrane
interactions into a unified activation state without consuming or destroying the driving input signals.
Furthermore, physical mechanisms for mechanical weight tuning, namely deformation limiting and intermediate-fluid
bulk-modulus regulation, provide a proposed physical medium for post-fabrication reconfigurability and future
neuroplasticity.

Mathematically, the proposed architecture maps directly onto the weighted-sum and sigmoidal activation forms of a
conventional artificial neural network under linear, constant-stiffness assumptions. When membrane stiffness varies
with deformation, the system generalises into a nonlinear, state-dependent neural computation element. While
physical constraints impose positivity and normalisation bounds on the reachable parameter space, the architecture
retains the expressiveness to serve as a building block for physical neural computation.

Future research directions include the physical fabrication and empirical characterisation of individual units to
analyse real-world fluid dynamics and mechanical response boundaries. Beyond single-unit mechanics, exploring the
integration of these neurons into multi-layered architectures presents a critical pathway toward fully autonomous,
self-learning mechanical neural networks capable of physical adaptation in soft robotic systems. Central to this
progression is establishing fully mechanical parameter-update mechanisms, enabling synaptic weights to adapt locally
without relying on external domain interventions.

\bibliographystyle{IEEEtran}
\bibliography{references}

% =====================================================================================================================
\clearpage
\begin{appendices}
\renewcommand{\theequation}{\thesection.\arabic{equation}}

\section{Intermediate-fluid bulk-modulus tuning}
\label{app:bulk}

\subsection{Definitions and sign convention}

The force-transfer-limiting mechanism consists of two flexible membranes separated by a confined intermediate fluid.
The input pressure $p_i$ acts on membrane~1, the activation-chamber pressure $p_a$ acts on membrane~2, and the
intermediate fluid develops the weight pressure $p_w$.

The effective areas of membranes 1 and 2 are denoted by $A_1$ and $A_2$, respectively. Their volumetric
deformations into the adjacent chambers are denoted by $\Delta V_1$ and $\Delta V_2$, with corresponding volumetric
stiffnesses $K_1$ and $K_2$. The deformation relevant to the activation response is the deformation of membrane~2,
\begin{equation}
  \Delta V_j = \Delta V_2 .
\end{equation}
The effective stiffness of the complete force-transfer mechanism is defined as
\begin{equation}
  K_j = \frac{F_j}{\Delta V_j},
\end{equation}
where $F_j$ is the effective force producing deformation into the activation chamber.

\subsection{Individual membrane force balances}

For membrane~1, the pressure difference between the input chamber and intermediate chamber produces the force
\begin{equation}
  F_1 = (p_i - p_w)\,A_1 .
\end{equation}
Using the volumetric stiffness relation $F_1 = K_1\,\Delta V_1$, the deformation of membrane~1 is
\begin{equation}
  \Delta V_1 = \frac{(p_i - p_w)\,A_1}{K_1} .
\end{equation}
For membrane~2, the pressure difference between the intermediate chamber and activation chamber produces the force
\begin{equation}
  F_2 = (p_w - p_a)\,A_2 .
\end{equation}
Similarly, $F_2 = K_2\,\Delta V_2$, and therefore
\begin{equation}
  \Delta V_j = \Delta V_2 = \frac{(p_w - p_a)\,A_2}{K_2} .
\end{equation}

\subsection{Intermediate-fluid compressibility}

Let $V_p$ denote the initial volume of the intermediate fluid and $B$ its bulk modulus. For small changes, the bulk
modulus relates the pressure change to the relative volume change according to
\begin{equation}
  \Delta p_w = -B\,\frac{\Delta V_p}{V_p} .
\end{equation}
Using the pressure increment relative to the initial weight pressure $p_{w,0}$,
\begin{equation}
  p_w - p_{w,0} = -B\,\frac{\Delta V_p}{V_p} .
\end{equation}
With the adopted deformation convention, the change in intermediate-fluid volume is
\begin{equation}
  \Delta V_p = \Delta V_2 - \Delta V_1 .
\end{equation}
Thus,
\begin{equation}
  p_w - p_{w,0} = -\frac{B}{V_p}\,(\Delta V_2 - \Delta V_1) .
\end{equation}
Substituting the individual membrane deformations gives
\begin{equation}
  p_w - p_{w,0} = -\frac{B}{V_p}\left[\frac{(p_w - p_a)\,A_2}{K_2} - \frac{(p_i - p_w)\,A_1}{K_1}\right] .
  \label{eq:weight-pressure}
\end{equation}
This equation determines the weight pressure $p_w$ as a function of the input pressure, activation pressure,
membrane properties, and intermediate-fluid bulk modulus.

\subsection{Effective stiffness}

The effective force associated with the complete two-membrane mechanism is written as
\begin{equation}
  F_j = F_1 + F_2 = (p_i - p_w)\,A_1 + (p_w - p_a)\,A_2 .
\end{equation}
The effective stiffness is therefore, with $\Delta V_j = \Delta V_2$,
\begin{equation}
  K_j = \frac{F_j}{\Delta V_j} = \frac{(p_i - p_w)\,A_1 + (p_w - p_a)\,A_2}{\Delta V_2} .
\end{equation}
Substituting the expression for $\Delta V_2$ gives
\begin{equation}
  K_j = \frac{(p_i - p_w)\,A_1 + (p_w - p_a)\,A_2}{\dfrac{(p_w - p_a)\,A_2}{K_2}}
      = K_2\,\frac{(p_i - p_w)\,A_1 + (p_w - p_a)\,A_2}{(p_w - p_a)\,A_2} .
\end{equation}
Because $p_w$ depends on the bulk modulus $B$, the effective stiffness is also a function of $B$, $K_j = K_j(B)$, and
the corresponding mechanical weight is
\begin{equation}
  W_j(B) = \frac{A_j}{K_j(B)} .
\end{equation}

\subsection{Limiting cases}

\subsubsection{Completely compressible intermediate fluid}
For a completely compressible intermediate fluid, $B = 0$. The intermediate fluid provides no resistance to
volumetric deformation and cannot sustain a pressure increase. Consequently, the force-transfer mechanism does not
produce a finite activation-side deformation associated with the input force, $\Delta V_j \rightarrow 0$. Therefore,
\begin{equation}
  K_j(B = 0) = \frac{F_j}{\Delta V_j} \rightarrow \infty .
\end{equation}
Thus, the completely compressible limit corresponds to an effectively infinitely stiff force-transfer mechanism.

\subsubsection{Completely incompressible intermediate fluid}
For a completely incompressible intermediate fluid, $B \rightarrow \infty$. The intermediate-fluid volume cannot
change, so $\Delta V_p \rightarrow 0$. Under the adopted sign convention, this requires the deformation of the two
membranes to be equal:
\begin{equation}
  \Delta V_1 = \Delta V_2 = \Delta V_j .
\end{equation}
The two membrane stiffnesses therefore act together in the fully coupled limit:
\begin{equation}
  F_j = K_1\,\Delta V_j + K_2\,\Delta V_j = (K_1 + K_2)\,\Delta V_j,
\end{equation}
and the effective stiffness becomes
\begin{equation}
  K_j(B \rightarrow \infty) = \frac{F_j}{\Delta V_j} = K_1 + K_2 .
\end{equation}
The limiting behaviour of the force-transfer mechanism is therefore
\begin{equation}
  K_j(0) \rightarrow \infty, \qquad K_j(\infty) = K_1 + K_2 .
\end{equation}
The intermediate-fluid bulk modulus consequently provides a tuneable parameter that controls the effective
stiffness of the two-membrane system and, through $W_j(B) = A_j/K_j(B)$, the corresponding mechanical weight. Bulk
modulus is conventionally defined as the ratio between pressure change and volumetric strain, with the sign
determined by the selected convention for compression and volume change.

% ---------------------------------------------------------------------------------------------------------------------
\section{Activation-chamber compliance}
\label{app:compliance}

Starting with the activation model with $\patm = 0$,
\begin{equation}
  p_a = \frac{\sum_j W_j p_j + W_b p_b}{\sum_j W_j + W_b + W_0},
\end{equation}
define $\pin = \sum_j W_j p_j + W_b p_b$, so that
\begin{equation}
  p_a = \left(\frac{1}{\sum_j W_j + W_b + W_0}\right)\pin = \alpha\,\pin,
  \qquad \alpha = \frac{1}{\sum_j W_j + W_b + W_0} .
\end{equation}
The activation-to-output function is
\begin{equation}
  \pout(p_a) = \plow + \frac{\phigh - \plow}{1 + \exp\bigl[-k\,(p_a - \pth)\bigr]} .
\end{equation}
Substituting $p_a = \alpha\,\pin$ yields
\begin{equation}
  \pout(\pin) = \plow + \frac{\phigh - \plow}{1 + \exp\bigl[-k\,(\alpha\,\pin - \pth)\bigr]} .
\end{equation}
The exponent can be rearranged as
\begin{equation}
  -k\,(\alpha\,\pin - \pth) = -k\alpha\left(\pin - \frac{\pth}{\alpha}\right) .
\end{equation}
Consequently,
\begin{equation}
  \pout(\pin) = \plow + \frac{\phigh - \plow}{1 + \exp\bigl[-k_{\mathrm{eff}}\,(\pin - p_{\mathrm{th,eff}})\bigr]},
\end{equation}
where
\begin{equation}
  k_{\mathrm{eff}} = k\,\alpha = \frac{k}{\sum_j W_j + W_b + W_0}
  \qquad\text{and}\qquad
  p_{\mathrm{th,eff}} = \frac{\pth}{\alpha} = \pth\Bigl(\sum_j W_j + W_b + W_0\Bigr) .
\end{equation}
The activation-chamber compliance therefore acts as a multiplicative scaling factor on the weighted input pressure.
Consequently, it changes the gradient of the sigmoid response: increasing the chamber compliance reduces the
effective gradient, while decreasing the compliance produces a steeper response.

% ---------------------------------------------------------------------------------------------------------------------
\section{Weight-membrane sensitivity}
\label{app:weight-sensitivity}

Define
\begin{equation}
  N = W_1 p_1 + \sum_{j\neq 1} W_j p_j + W_b p_b + W_0 \patm,
  \qquad
  D = W_1 + \sum_{j\neq 1} W_j + W_b + W_0,
\end{equation}
so that the activation pressure is $p_a = N/D$. Differentiating with respect to $W_1$, while treating all other
pressures and weights as constant, gives
\begin{equation}
  \frac{\partial p_a}{\partial W_1}
    = \frac{\dfrac{\partial N}{\partial W_1}\,D - N\,\dfrac{\partial D}{\partial W_1}}{D^2} .
\end{equation}
Since $\partial N/\partial W_1 = p_1$ and $\partial D/\partial W_1 = 1$, it follows that
\begin{equation}
  \frac{\partial p_a}{\partial W_1} = \frac{p_1 D - N}{D^2} .
\end{equation}
Using $N = p_a D$,
\begin{equation}
  \frac{\partial p_a}{\partial W_1} = \frac{p_1 D - p_a D}{D^2},
\end{equation}
which simplifies to
\begin{equation}
  \boxed{\frac{\partial p_a}{\partial W_1} = \frac{p_1 - p_a}{D}} .
\end{equation}
The mechanical weight is related to the membrane stiffness by $W_1 = A_1/K_1$. Therefore,
\begin{equation}
  \frac{\partial W_1}{\partial K_1} = -\frac{A_1}{K_1^2} .
\end{equation}
Applying the chain rule,
\begin{equation}
  \frac{\partial p_a}{\partial K_1} = \frac{\partial p_a}{\partial W_1}\,\frac{\partial W_1}{\partial K_1}
    = -\frac{A_1}{K_1^2}\,\frac{p_1 - p_a}{D} .
\end{equation}
Hence,
\begin{equation}
  \boxed{\frac{\partial p_a}{\partial K_1}
    = -\frac{A_1\,(p_1 - p_a)}{K_1^2\bigl(W_1 + \sum_{j\neq 1} W_j + W_b + W_0\bigr)}} .
\end{equation}

% ---------------------------------------------------------------------------------------------------------------------
\section{Input-pressure sensitivity}
\label{app:input-sensitivity}

With $\patm = 0$, the activation pressure is
\begin{equation}
  p_a = \frac{\sum_j W_j p_j + W_b p_b}{\sum_j W_j + W_b + W_0} .
\end{equation}
The denominator is independent of the input pressures $p_i$. Differentiating with respect to $p_i$ therefore gives
\begin{equation}
  \frac{\partial p_a}{\partial p_i}
    = \frac{\partial\bigl(\sum_j W_j p_j + W_b p_b\bigr)/\partial p_i}{\sum_j W_j + W_b + W_0} .
\end{equation}
Since
\begin{equation}
  \frac{\partial}{\partial p_i}\Bigl(\sum_j W_j p_j\Bigr) = W_i
\end{equation}
and the bias term $W_b p_b$ is independent of $p_i$,
\begin{equation}
  \frac{\partial p_a}{\partial p_i} = \frac{W_i}{\sum_j W_j + W_b + W_0} .
\end{equation}
Using $\alpha$,
\begin{equation}
  \boxed{\frac{\partial p_a}{\partial p_i} = \alpha\,W_i} .
\end{equation}

% ---------------------------------------------------------------------------------------------------------------------
\section{Equivalent ANN weights and bias}
\label{app:ann}

With $\patm = 0$, the activation pressure is
\begin{equation}
  p_a = \alpha\Bigl(\sum_j W_j p_j + W_b p_b\Bigr),
\end{equation}
where $\alpha$ is defined in \cref{eq:alpha}. Differentiating with respect to an input pressure $p_j$ gives
\begin{equation}
  \frac{\partial p_a}{\partial p_j} = \alpha\,W_j .
\end{equation}
The activation pressure can therefore be expressed in the artificial-neural-network form
\begin{equation}
  p_a = \sum_j \left(\frac{\partial p_a}{\partial p_j}\right) p_j + b .
\end{equation}
Substituting the derived input sensitivities gives
\begin{equation}
  p_a = \sum_j \alpha\,W_j\,p_j + b = \alpha \sum_j W_j\,p_j + b .
\end{equation}
Equating this expression with the original activation-pressure equation gives
\begin{equation}
  \alpha \sum_j W_j\,p_j + b = \alpha\Bigl(\sum_j W_j\,p_j + W_b\,p_b\Bigr) .
\end{equation}
After cancelling the common input-pressure term,
\begin{equation}
  \boxed{b = \alpha\,W_b\,p_b = \frac{W_b\,p_b}{\sum_j W_j + W_b + W_0}} .
\end{equation}
Consequently, the activation pressure has the equivalent artificial-neural-network form
\begin{equation}
  \boxed{p_a = \sum_j w_j\,p_j + b},
\end{equation}
where the effective ANN weights are
\begin{equation}
  \boxed{w_j = \frac{\partial p_a}{\partial p_j} = \alpha\,W_j = \frac{W_j}{\sum_k W_k + W_b + W_0}} .
\end{equation}
Applying the activation-to-output function then gives
\begin{equation}
  \boxed{\pout = F(p_a) = F\Bigl(\sum_j w_j\,p_j + b\Bigr)} .
\end{equation}

% ---------------------------------------------------------------------------------------------------------------------
\section{Equal-pressure hyperplane}
\label{app:hyperplane}

Consider the pre-activation relation
\begin{equation}
  p_a = \sum_j w_j\,p_j + w_b\,p_b + w_0\,p_0
\end{equation}
subject to
\begin{equation}
  \sum_j w_j + w_b + w_0 = 1 .
\end{equation}
Let all input and reference pressures be equal to some arbitrary pressure $p$, $p_j = p_b = p_0 = p$. Substituting
into the pre-activation equation gives
\begin{equation}
  p_a = \sum_j w_j\,p + w_b\,p + w_0\,p = \Bigl(\sum_j w_j + w_b + w_0\Bigr)p .
\end{equation}
Using the normalisation constraint, $p_a = p$. Therefore
\begin{equation}
  \boxed{p_1 = p_2 = \cdots = p_N = p_b = p_0 = p_a}
\end{equation}
satisfies the pre-activation equation for any value of $p$. Hence, the hyperplane necessarily contains this
one-dimensional equal-pressure diagonal.

Thus, the normalisation constraint does not merely restrict the magnitudes of the weights; it geometrically
constrains the family of representable hyperplanes to those passing through the equal-pressure diagonal. For a
fixed value of $p_b$, the hyperplane is constrained to pass through the corresponding point on the equal-pressure
diagonal,
\begin{equation}
  (p_1, p_2, \ldots, p_N, p_a) = (p_b, p_b, \ldots, p_b, p_b) .
\end{equation}
Since $p_b$ is a mechanically tuneable parameter, this point can be moved along the diagonal, providing control over
the position of the hyperplane. The input weights $w_j$, meanwhile, determine its orientation, with $w_b$ determined
by the normalisation constraint. Thus, the mechanical parameters provide control over both the orientation and
position of the hyperplane, although only within the restrictions imposed by the non-negative mechanical parameters
and the normalisation constraint.

\end{appendices}
\end{document}
